\clearpage
\onecolumn
\thispagestyle{empty}
\vspace*{\fill}
\begin{center}
	{\color{radblue}\rule{\linewidth}{2pt}}\par
	\vspace{6mm}
	{\Huge\bfseries\color{radblue} Consideraciones Generales}\par
	\vspace{6mm}
	{\color{radblue}\rule{\linewidth}{2pt}}
\end{center}
\vspace*{\fill}
\clearpage
\twocolumn
\pagestyle{contenido}%%%%%%%%%%%%%%%%%%%%%%%%%%%%%%%%%%%%%%%%%%%%%%%%%%%%%%%%%%%




%  SECCIÓN: Consideraciones Generales
%  Archivo: secciones/consideraciones_generales.tex
%  Sin preámbulo ni \begin{document} — lo maneja main.tex
	%%%%%%%%%%%%%%%%%%%%%%%%%%%%%%%%%%%%%%%%%%%%%%%%%%%%%%%%%%%
	
	\section{Consideraciones Generales}
	
	%----------------------------------------------------------
	\subsection{Dato Clínico y Enfoques Complementarios}
	%----------------------------------------------------------
	
	Una vez realizado el par radiológico se pueden realizar enfoques complementarios, los cuales sirven para despejar dudas y aportar información diferente. El dato clínico determina cuál es el enfoque complementario correspondiente en caso de ser necesario.
	
	\textbf{Datos clínicos más comunes:}
	\begin{itemize}
		\item Traumatismo
		\item Fractura
		\item Luxación
		\item Dolor
		\item Derrame articular
		\item Osteoporosis
		\item Gota
		\item Metástasis
		\item Tumores
	\end{itemize}
	
	%----------------------------------------------------------
	\subsection{Valor Diagnóstico}
	%----------------------------------------------------------
	
	Se deben cumplir ciertos requisitos para otorgarle valor diagnóstico a la radiografía:
	\begin{itemize}
		\item Marca de la derecha
		\item Región comprendida
		\item Posición lograda
		\item Técnica adecuada
	\end{itemize}
	
	El valor diagnóstico de la radiografía siempre está relacionado al dato clínico.
	
	\textbf{Importante:} Una Rx cortada no tiene valor diagnóstico ya que no cumple con la región comprendida, pero siempre se debe relacionar la Rx con su dato clínico para determinar si se le otorga o no valor.
	
	%----------------------------------------------------------
	\subsection{Parámetros de Exposición en Radiografía Convencional}
	%----------------------------------------------------------
	
	\textbf{kVp (kilovoltaje pico):} determina la calidad del haz de rayos X, es decir, su poder de penetración. Un aumento del kVp produce fotones de mayor energía, lo que amplía la escala de grises y disminuye el contraste. Valores bajos de kVp generan una escala de grises corta, resultando en alto contraste (``escala corta''). El kVp influye además en la cantidad de radiación dispersa producida: a mayor kVp, mayor dispersión, lo que deteriora la calidad de imagen si no se adoptan medidas correctoras.
	
	\textbf{mAs (carga eléctrica):} gobierna la cantidad de radiación, es decir, la dosis que recibe el paciente. El uso de alto kVp permite reducir los mAs necesarios para un ennegrecimiento adecuado, disminuyendo la dosis al paciente; mientras que el uso de bajo kVp obliga a incrementar los mAs para compensar la menor penetración del haz, con el consiguiente aumento de dosis.
	
	\textbf{Bucky (rejilla antidifusora):} se emplea cuando el espesor de la región anatómica genera cantidad significativa de radiación dispersa. Al eliminar los fotones dispersos, mejora y aumenta el contraste. Sin embargo, absorbe también una fracción de la radiación primaria útil, exigiendo mayor dosis de exposición. Además, introduce una leve reducción de la definición espacial por la trama de sus láminas. Su uso solo está justificado cuando el espesor de la parte explorada realmente lo amerita.
	
	\textbf{Correcciones en la práctica clínica:}
	\begin{itemize}
		\item \textbf{Imagen sobreexpuesta:} se corrige reduciendo los mAs o, si el contraste lo permite, aumentando el kVp.
		\item \textbf{Imagen subexpuesta:} se corrige principalmente aumentando los mAs, ya que la cantidad de radiación que llegó al detector fue insuficiente.
	\end{itemize}
	
	%----------------------------------------------------------
	\subsection{Marcación Radiológica}
	%----------------------------------------------------------
	
	La marca radiológica identifica el lado del paciente. Como regla universal, al montar la imagen en el negatoscopio siguiendo la convención estándar, la marca queda invariablemente a la \textbf{izquierda del observador}, independientemente de la proyección utilizada.
	
	\subsubsection*{Altura}
	
	La marca se coloca \textbf{arriba} cuando el paciente está en bipedestación o a 45°, y \textbf{abajo} cuando está sentado o en decúbito supino.
	
	\subsubsection*{Perfiles}
	
	El perfil derecho se marca \textbf{anterior} y el perfil izquierdo se marca \textbf{posterior}. En el par radiológico solo se marca el frente.
	
	\subsubsection*{Particularidades Anatómicas}
	
	\begin{itemize}
		\item \textbf{Mano (PA):} marca del lado del pulgar para la mano izquierda; del lado del meñique para la mano derecha.
		\item \textbf{Antebrazo (AP):} marca del lado del radio (externo) para el derecho; del lado del cúbito (interno) para el izquierdo.
		\item \textbf{Pierna (AP):} marca del lado del peroné (externo) para el derecho; del lado de la tibia (interno) para el izquierdo.
		\item \textbf{Hombro (AP):} marca en el lado externo para el derecho; en el lado interno para el izquierdo.
		\item \textbf{Tórax (PA):} marca del lado contrario a la silueta cardíaca.
	\end{itemize}
	
	%----------------------------------------------------------
	\subsection{Procedimientos Generales}
	%----------------------------------------------------------
	
	Previo a llamar al paciente, se debe alinear el tubo con el Bucky, luego con el paciente y finalmente el rayo central (RC) en el punto de incidencia correspondiente. Este procedimiento aplica tanto para el Bucky de pared como para la mesa.
	
	\begin{enumerate}
		\item Llamar al paciente por nombre y apellido, verificando su identidad junto con la cédula de identidad.
		\item Hacerlo pasar, saludarlo y presentarse, dirigiéndose de forma cordial y profesional.
		\item Verificar que el pedido sea correcto: confirmar la región a estudiar y, en caso de que no especifique lateralidad, consultársela al paciente. Se puede indagar brevemente sobre el motivo de consulta.
		\item Medir el espesor de la región a explorar.
		\item Mientras el paciente se cambia, calcular la técnica radiográfica.
		\item Posicionar al paciente y darle las indicaciones previas a la exposición: no moverse, mantener la posición durante la apnea o respiración indicada, entre otras.
		\item Cerrar la bandeja del chasis.
		\item Cerrar las puertas del gabinete.
		\item Realizadas las exposiciones, indicar al paciente que puede relajarse y acomodarse según corresponda al estudio.
	\end{enumerate}
	
	\subsubsection*{Pregunta de Embarazo}
	
	A toda paciente de sexo femenino en edad fértil, comprendida aproximadamente entre los 12 y 50 años, se le debe consultar por protocolo si existe posibilidad de embarazo. La pregunta se formula de forma discreta: \textit{``Disculpe, es por protocolo: ¿hay alguna posibilidad de embarazo?''}
	
	En caso de respuesta afirmativa, se le consultará si el médico solicitante fue informado y si le explicó los riesgos asociados al estudio. Cuando la paciente sea menor de edad y se encuentre acompañada por un familiar, se solicitará amablemente que el acompañante se retire por unos minutos antes de realizar la consulta.
	
	\subsubsection*{Preparación del Paciente}
	
	Se solicitará al paciente que retire todos los elementos que puedan generar artefactos en la imagen:
	\begin{itemize}
		\item Anillos y relojes. De no poder retirarlos, el paciente deberá lavarse con agua y jabón.
		\item Prendas con elementos metálicos: chapas, brillos, lentejuelas y estampados metálicos.
		\item Corpiños con aro, verificando siempre que el paciente haya cumplido con la indicación.
		\item Pantalones de jean, dado que la cremallera genera artefactos significativos. Los pantalones deportivos o de tela pueden mantenerse puestos.
	\end{itemize}
	
	\subsubsection*{Consideraciones Prácticas}
	
	\begin{itemize}
		\item \textbf{Paciente con dolor:} mantenerlo tranquilo, hablarle con calma e informarle que el estudio es rápido y necesario para el diagnóstico.
		\item \textbf{Paciente pediátrico acompañado por madre embarazada:} la madre debe retirarse. Ingresará otro acompañante, a quien se le colocará el chaleco plomado para permanecer durante la exposición. Es frecuente que los niños lloren durante el procedimiento; forma parte del manejo habitual.
		\item \textbf{Limpieza del chasis y la mesa radiográfica:} cuando se haya apoyado material potencialmente contaminado (pie diabético, lesiones infectadas, entre otros) se debe limpiar el chasis y la mesa con algodón embebido en alcohol. No se debe verter alcohol directamente sobre el chasis. En lo posible, limpiar siempre delante del paciente.
		\item \textbf{Uso de guantes:} si se utilizan guantes durante la atención al paciente, deben retirarse antes de manipular el tubo de rayos X.
	\end{itemize}
	
	%----------------------------------------------------------
	\subsection{Leyes de Proyección}
	%----------------------------------------------------------
	
	Las leyes de proyección describen el comportamiento geométrico del haz de rayos X al interactuar con el objeto y el receptor de imagen, condicionando la fidelidad y las limitaciones inherentes a toda imagen radiológica.
	
	\begin{itemize}
		\item \textbf{Proyección cónica:} los rayos X se originan en un foco puntual y se expanden en forma de abanico. La imagen resultante es siempre mayor que el objeto real, fenómeno conocido como \textbf{magnificación}, aunque en condiciones estándar esta es mínima.
		
		\item \textbf{Superposición y sumación:} la imagen radiológica proyecta un objeto tridimensional sobre un plano bidimensional. Todas las estructuras atravesadas por el RC se aplanan en un único plano, generando superposición de estructuras y sumación de densidades. Ejemplo representativo: superposición de los cuerpos vertebrales con el contenido toracoabdominal y mediastínico.
		
		\item \textbf{Ley de las tangencias:} el contorno de una estructura solo se visualiza con nitidez cuando el RC incide de forma tangente a su superficie. Esta ley se cumple siempre. En la proyección de pelvis, la rotación interna de los fémures permite que el RC pase tangencial al cuello femoral, visualizando sus bordes óseos con precisión. En la columna panorámica, se cumple en los bordes de los cuerpos vertebrales.
		
		\item \textbf{Ley de conservación de la forma:} para evitar la distorsión geométrica, el RC debe mantener una relación ortogonal (a 90°) con el receptor de imagen. Cuando esta condición se cumple, la forma del objeto se conserva de manera fiel, con la única limitación de una magnificación mínima inherente a la proyección cónica.
		
		\item \textbf{Ángulo de incidencia:} se aplica cuando se requiere elongar una estructura o desplazar una superposición mediante la angulación deliberada del RC. En condiciones estándar, donde el RC incide perpendicular al receptor, esta ley no se aplica y se prioriza la conservación de la forma anatómica.
	\end{itemize}
	
	%----------------------------------------------------------
	\subsection{Examen: Enfoques y DC Frecuentes}
	%----------------------------------------------------------
	
	Al ingresar al examen, las tres imágenes ya se encuentran colocadas y el DC es comunicado de forma oral por el evaluador. Se recomienda seguir estrictamente los puntos indicados en la rúbrica para evitar preguntas complementarias.
	
	\textbf{Enfoques y diagnósticos clínicos relevados como frecuentes en instancias previas:}
	\begin{itemize}
		\item Frente de tórax con catéter venoso central
		\item Hueso nasal perfil
		\item Vértices pulmonares (DC: tuberculosis)
		\item Perfil de sacro-cóccix
		\item Maxilar inferior desenfilado
		\item Perfil de tórax
		\item Perfil de silla turca
		\item Frente de columna cervical (mal posicionado, DC: traumatismo)
		\item Outlet
		\item Oblicua de vejiga
		\item Caldwell modificado
		\item Frente de lumbosacra (tornillos)
		\item Columna lumbosacra perfil
		\item Frente de cráneo (mal posicionado)
		\item Funcional de lumbar en flexión (hiperflexión)
		\item Abdomen de pie (DC: oclusión intestinal)
		\item Waters (DC: traumatismo)
		\item Pelvis AP común (prótesis derecha)
		\item Cigomático unilateral
		\item Hirtz modificado para cigomático
		\item Frente de sacro
		\item Caldwell a 15°
		\item Frente de abdomen
		\item Pie (DC: traumatismo)
		\item Perfil de columna cervical
		\item Waters-Mahoney
		\item Abdomen en decúbito
		\item Panorámica
		\item Schüller
		\item Omalgia (ambas rotaciones)
		\item Desfiladero
		\item Alar y obturatriz
	\end{itemize}
	
	\textbf{Importante:} Si te preguntan si se puede mejorar, siempre responder que sí, ya que siempre algo en la Rx se puede mejorar.
	
\clearpage